%----------------------------------------------------------------------------------------
% VOORKANT: dezelfde indeling voor alle handboeken. Elke cursus leest dit bestand rechtstreeks
% (via Opmaak/opmaak.tex): pas het dus alleen hier aan.
%
% Per onderwerp:  motieven/<naam>.tex  kunstwerk, bovenrand van de band en vormen
% Per cursus:     voorkant-gegevens.tex  rubriek, richtingen, graad, uren, titel en \motiefnaam
% Coördinaten in cm, met de oorsprong linksonder op de pagina (A4: 21 x 29,7).
%----------------------------------------------------------------------------------------
\providecommand{\omslagsamples}{120}

% Hoogte van de bovenrand op plaats #1, in macro #2. Handig om iets op de rand te zetten.
\newcommand{\randwaarde}[2]{\def\x{#1}\pgfmathsetmacro{#2}{\omslagrand}}
\newcommand{\randlijn}{%
  \draw[white,line width=2pt,line cap=round,line join=round]
    plot[domain=0:21,samples=\omslagsamples,variable=\x] (\x,{\omslagrand});}
\newcommand{\standaardband}{%
  \fill[accent] (0,0) -- plot[domain=0:21,samples=\omslagsamples,variable=\x] (\x,{\omslagrand}) -- (21,0) -- cycle;}

% Witte lijnen met een rand in de bandkleur, zodat ze ook op een licht schilderij zichtbaar blijven.
\tikzset{
  omslaglijn/.style={draw=white,line width=#1,preaction={draw=accent,line width=#1+1.8pt}},
  omslaglijn/.default=1pt,
  omslagstip/.style={fill=white,draw=accent,line width=.7pt},
  omslaglabel/.style={fill=accent,text=white,inner sep=2pt,rounded corners=1.5pt,font=\small},
}

\ifdefined\motiefnaam\InputIfFileExists{\hoofdmap Ontwerp/voorkanten/motieven/\motiefnaam}{}{}\fi

\newcommand{\maakomslag}{%
  \begingroup\thispagestyle{empty}%
  % de kleur van de graad (voorkant-gegevens.tex); zonder \omslagaccent blijft de kleur van de 3de graad
  \ifdefined\omslagaccent\definecolor{accent}{HTML}{\omslagaccent}\definecolor{fel}{HTML}{\omslagfel}\fi
  \begin{tikzpicture}[remember picture,overlay]
  \begin{scope}[shift={(current page.south west)}]
    \node[anchor=north west,inner sep=0pt] at (0,29.7)
      {\includegraphics[width=\paperwidth]{\omslagbestand}};
    \ifcsname omslagband\endcsname\omslagband\else\standaardband\fi
    \omslagmotief
    % titel, met de rubriek erboven en de versie eronder
    \node[name=omslagtitelnode,anchor=south west,inner sep=0pt,text=white,align=left,
      font=\kopletter\fontsize{\omslagtitelgrootte}{\numexpr\omslagtitelgrootte*107/100\relax}\selectfont]
      at (2,5.2) {\strut\omslagtitel\strut};
    \node[anchor=south west,inner sep=0pt,text=fel,font=\sffamily\bfseries\normalsize]
      at ([yshift=2mm]omslagtitelnode.north west)
      {Wiskunde\ifdefempty{\cursusrubriek}{}{ \textbullet\ \cursusrubriek}};
    \node[anchor=north west,fill=fel,text=inkt,rounded corners=2.5pt,inner xsep=2.5mm,inner ysep=1.2mm,
      font=\kopletter\large] at (2,5.15) {\cursusondertitel};
    % de gegevens
    \node[anchor=south west,inner sep=0pt,text=white] at (2,1.45)
      {\sffamily\small\renewcommand{\arraystretch}{1.15}\hyphenpenalty=10000\exhyphenpenalty=10000%
       \begin{tabular}{@{}>{\bfseries}l@{\hspace{4mm}}>{\RaggedRight\arraybackslash}p{9.2cm}@{}}
         Richtingen & \cursusrichtingen \\
         Graad      & \cursusgraad \\
         Wiskunde   & \cursusuren
       \end{tabular}};
    \node[anchor=south west,inner sep=0pt,text=white,font=\sffamily\tiny] at (2,.75) {Omslag: \omslagbron};
    \node[anchor=south east,inner sep=0pt] at (19,1.5)
      {\includegraphics[width=4.6cm]{axiomath-logo-wit.pdf}};
  \end{scope}
  \end{tikzpicture}%
  \vfill\endgroup\clearpage}
